%

\subsection{Schweikart i Taurinus} % lang-pl
\label{sekcja_schweikart_taurinus}

Opowiemy o dwóch niemieckich prawnikach, Ferdinandzie Karlu Schweikarcie i jego bratanku, Franzie Adolphie Taurinusie.

Schweikart w 1807 roku opublikuje pracę ,,Die Theorie der Parallellinien, nebst dem Vorschlage ihrer Verbannung aus der Geometrie''\footnote{Teoria linii równoległych wraz z propozycją usunięcia ich z geometrii.}.
W 1818 roku przekaże Gaussowi za pośrednictwem swojego studenta, Christiana Gerlinga, krótką notatkę o~nowej geometrii, zwanej przez niego \emph{geometrią astralną}, w której suma kątów trójkąta jest mniejsza niż $\pi$.
\index[persons]{Schweikart, Ferdinand Karl}%
\index[persons]{Gerling, Christian Ludwig}%
Taurinus napisze w 1824 roku do Gaussa, który zachęci go do dalszych badań, ale poprosi o dyskrecję.
Taurinus opublikuje w 1825 roku \emph{Theorie der Parallellinien}, a w 1826 roku \emph{Geometriae prima elementa}.
Podejmie domysł Lamberta: zbada geometrię na sferze o urojonym promieniu, nazwaną przez siebie logarytmiczno-sferyczną, i wyprowadzi w niej między innymi hiperboliczne twierdzenie cosinusów.
Sam jednak pozostanie przy przekonaniu, że geometria euklidesowa jest wyjątkowa, a rozczarowany brakiem uznania zniszczy większość egzemplarzy swojej drugiej książki.
\index[persons]{Taurinus, Franz Adolph}%

% Odcinek, którego kąt równoległości wynosi $\pi/4$, nazywa się dziś odcinkiem Schweikarta \cite[s. 211]{greenberg_2010}.

% https://en.wikipedia.org/wiki/Ferdinand_Karl_Schweikart
% https://en.wikipedia.org/wiki/Franz_Taurinus

%
